\chapter{The pub/sub library}
\label{ch:library}

A low-latency, multi-threaded, event-driven application framework for C++17,
built around the reactor pattern. It provides inter-thread communication,
inter-process communication, pub/sub messaging, timers, high availability, and a
binary serialisation DSL --- all designed for environments where heap allocation
on the hot path is not acceptable.

The rest of this chapter is not yet written, except for the placeholders below.

\section{Pub/sub}
\label{sec:pubsub}

\subsection{Overview}
pubsub\_itc\_fw provides a topic-based publish/subscribe messaging facility as one of the core services
of the framework, alongside inter-thread communication (ITC), inter-process communication (IPC), and timers.
The pub/sub facility allows a component (a publisher) to emit a message tagged with a topic,
without needing to know which components (subscribers), if any, are interested in that topic.
Subscribers register interest in one or more topics and receive a copy of every message published
to those topics.

This decoupling is the central value of the feature: publishers and subscribers do not hold
direct references to one another, do not need to be co-located in the same thread or process,
and do not need to be started or stopped in any particular order relative to each other.
A component's only coupling to the rest of the system is the set of topic names it produces or consumes.

Because pub/sub sits on top of the same reactor event loop as the other primitives,
a component that publishes or subscribes to topics does so using the same non-blocking,
single-threaded-per-reactor programming model used everywhere else in the framework — there is
no separate thread pool or blocking call dedicated to messaging.

At a functional level, the pub/sub facility is expected to provide:

\begin{itemize}
\item Topics — named channels that categorise messages. A topic is a logical concept, not a physical connection; any number of publishers and subscribers may use the same topic name.

\item Publish — a component sends a message tagged with a topic name. The publisher does not block waiting for delivery and does not receive any acknowledgement that a specific subscriber processed the message.
Subscribe — a component registers interest in a topic (or set of topics) and subsequently receives every message published to that topic for as long as the subscription is active.
\item Delivery — the implementation delivers messages via unicast fanout: the framework sends an individual copy of each message to each interested subscriber over its own connection, rather than using a network-level multicast/broadcast mechanism. This keeps delivery reliable and simple to reason about at the cost of publisher-side work scaling with subscriber count.

\item Message ordering — messages published to a given topic are delivered to a given subscriber in the order they were published.

\item Serialisation — message payloads are defined using the framework's binary serialisation DSL and compiled to C++17 encode/decode code, rather than hand-written wire formats. This keeps the pub/sub payload format consistent with the one used for ITC and IPC messages elsewhere in the framework.
\end{itemize}

The remainder of this section specifies the pub/sub facility in detail:

\begin{itemize}
\item Topic naming and lifecycle
\item Publisher and subscriber APIs
\item Subscription management (add/remove, late-joining subscribers)
\item Delivery semantics and failure handling (slow subscribers, disconnects)
\end{itemize}

\begin{description}
  \item[What it is.] A publisher exposes named topics; subscribers connect, replay the history
    they have missed, and then stream new records live. What a topic is, what a record is, and
    what a subscriber is entitled to expect.
  \item[Why the fanout is unicast.] The reference system this venue's semantics are drawn from
    uses reliable multicast over the user datagram protocol, together with a broker. This venue
    copies the semantics and not the mechanism. The reasoning is to be written out in full here,
    because at present it survives only as a parenthesis in a superseded working note.
  \item[The alternatives, and why each was rejected.] Each rejection is to name the ground it
    rests on, and that ground is to be one a reader can check. A rejection that rests on a
    mistaken premise discredits the ones that do not.
  \item[An outline of the functionality.] Delivery and ordering, durability and replay, fanout
    and backpressure, failover, and subscriber lifecycle and identity.
\end{description}
